
\subsection{The affine $\mathfrak{sl}_2$ ordered bar complex
and $Y_\hbar(\mathfrak{sl}_2)$}
\label{subsec:sl2-yangian}

We now unwind the ordered bar complex for the affine Kac--Moody algebra
$\widehat{\mathfrak{sl}}_2$ at level~$k$, producing the Yangian
$Y_\hbar(\mathfrak{sl}_2)$ as the open-colour Koszul dual. This is
the simplest non-abelian instance and already exhibits every essential
feature: the collision residue, the Yang $R$-matrix, the RTT
presentation, the Gauss decomposition, the PBW basis, and evaluation
modules. We emphasize that $\widehat{\mathfrak{sl}}_2$ is an
$E_\infty$-chiral algebra (a local vertex algebra in the sense of
Beilinson--Drinfeld); the Yangian is the $E_1$-chiral Koszul dual
, the \emph{output} of applying $E_1$-chiral Koszul duality to an
$E_\infty$ input.

Throughout this subsection, $\mathfrak g=\mathfrak{sl}_2$ with standard
basis $\{e,f,h\}$, structure constants $[e,f]=h$, $[h,e]=2e$,
$[h,f]=-2f$, Killing form $\kappa(x,y)=\tfrac12\operatorname{tr}(\operatorname{ad}(x)\operatorname{ad}(y))$
normalised so that
$\kappa(e,f)=\kappa(f,e)=1$, $\kappa(h,h)=2$. The dual
Coxeter number is $h^\vee=2$.

\subsubsection*{The OPE and its collision residue}

The affine $\widehat{\mathfrak{sl}}_2$ OPE at level $k$ is
\begin{equation}\label{eq:sl2-ope}
J^a(z)\,J^b(w)\;\sim\;
\frac{k\,\kappa^{ab}}{(z-w)^2}\;+\;\frac{f^{ab}{}_c\,J^c(w)}{z-w}\,.
\end{equation}
By the $d\log$ absorption principle (the bar complex uses $d\log(z{-}w)$
propagators, not $dz$ propagators; the kernel absorbs one power of the
pole), the ordered bar complex uses
the collision residue $r^{\mathrm{coll}}(z)$. The OPE~\eqref{eq:sl2-ope}
has a double pole (from the level) and a simple pole (from the Lie
bracket); after absorption, the collision residue retains a simple pole
from the double-pole term and a regular term from the simple-pole term:
\begin{equation}\label{eq:sl2-collision-residue}
r^{\mathrm{coll}}(z)\;=\;\frac{\hbar\,\Omega}{z}\;+\;\text{(regular)},
\end{equation}
where
\begin{equation}\label{eq:sl2-casimir}
\Omega\;=\;e\otimes f\;+\;f\otimes e\;+\;\tfrac12\,h\otimes h
\;\in\;\mathfrak g\otimes\mathfrak g
\end{equation}
is the Casimir element (the inverse of the Killing form), and
\begin{equation}\label{eq:sl2-hbar}
\hbar\;=\;\frac{1}{k+h^\vee}\;=\;\frac{1}{k+2}\,.
\end{equation}
The regular terms at order $z^0$ reproduce the Lie bracket
$[x,y]=f^{ab}{}_c$.

\subsubsection*{Degree 1: generators}

The ordered bar complex $\Barch^{\mathrm{ord}}(\widehat{\mathfrak{sl}}_2)$
in degree~$1$ is spanned by
\[
\susp^{-1}e,\qquad \susp^{-1}f,\qquad \susp^{-1}h,
\]
where $\susp^{-1}$ denotes the desuspension (bar degree shift). Since
$\Barch^{\mathrm{ord}}_0=\KK$ and the augmentation ideal starts in
conformal weight~$0$, the bar differential on degree-$1$ elements
vanishes:
\[
d(\susp^{-1}x)\;=\;0\qquad\text{for all }x\in\{e,f,h\}.
\]
The three classes $[\susp^{-1}e]$, $[\susp^{-1}f]$, $[\susp^{-1}h]$ in
$H^1(\Barch^{\mathrm{ord}})$ are the degree-$1$ bar cohomology generators.
Under the identification with the Yangian, they correspond to the
leading modes of the Drinfeld generating functions:
\[
E(u)\;=\;\sum_{r\ge 0}E^{(r)}\,u^{-r-1},\qquad
F(u)\;=\;\sum_{r\ge 0}F^{(r)}\,u^{-r-1},\qquad
H(u)\;=\;1+\sum_{r\ge 0}H^{(r)}\,u^{-r-1},
\]
where $E^{(0)}\leftrightarrow [\susp^{-1}e]$,
$F^{(0)}\leftrightarrow [\susp^{-1}f]$,
$H^{(0)}\leftrightarrow [\susp^{-1}h]$, and the spectral parameter
$u=1/z$ is the inverse of the curve coordinate.

\subsubsection*{Degree 2: the ordered bar differential and Lie relations}

\begin{computation}[Degree-$2$ ordered bar complex of
$\widehat{\mathfrak{sl}}_2$; \ClaimStatusProvedHere]
\label{comp:ordered-bar-sl2}
In degree~$2$, the ordered bar complex has $3\times 3=9$ generators:
\[
\susp^{-1}x\otimes \susp^{-1}y,\qquad x,y\in\{e,f,h\}.
\]
The ordered complex keeps $\susp^{-1}e\otimes \susp^{-1}h$ and
$\susp^{-1}h\otimes \susp^{-1}e$ as \textbf{distinct} elements: this is
the essential difference from the symmetric bar complex, which would
identify these up to sign.

The bar differential $d:\Barch^{\mathrm{ord}}_2\to\Barch^{\mathrm{ord}}_1$
extracts the zeroth product (Lie bracket) from the OPE:
\begin{align}
d(\susp^{-1}e\otimes \susp^{-1}f)
&\;=\;-\susp^{-1}[e,f]\;=\;-\susp^{-1}h,
\label{eq:d-ef}\\
d(\susp^{-1}f\otimes \susp^{-1}e)
&\;=\;-\susp^{-1}[f,e]\;=\;+\susp^{-1}h,
\label{eq:d-fe}\\
d(\susp^{-1}h\otimes \susp^{-1}e)
&\;=\;-\susp^{-1}[h,e]\;=\;-2\,\susp^{-1}e,
\label{eq:d-he}\\
d(\susp^{-1}e\otimes \susp^{-1}h)
&\;=\;-\susp^{-1}[e,h]\;=\;+2\,\susp^{-1}e,
\label{eq:d-eh}\\
d(\susp^{-1}h\otimes \susp^{-1}f)
&\;=\;-\susp^{-1}[h,f]\;=\;+2\,\susp^{-1}f,
\label{eq:d-hf}\\
d(\susp^{-1}f\otimes \susp^{-1}h)
&\;=\;-\susp^{-1}[f,h]\;=\;-2\,\susp^{-1}f,
\label{eq:d-fh}\\
d(\susp^{-1}e\otimes \susp^{-1}e)
&\;=\;-\susp^{-1}[e,e]\;=\;0,
\label{eq:d-ee}\\
d(\susp^{-1}f\otimes \susp^{-1}f)
&\;=\;-\susp^{-1}[f,f]\;=\;0,
\label{eq:d-ff}\\
d(\susp^{-1}h\otimes \susp^{-1}h)
&\;=\;-\susp^{-1}[h,h]\;=\;0.
\label{eq:d-hh}
\end{align}
The sign convention is $d(\susp^{-1}x\otimes \susp^{-1}y)=-\susp^{-1}(x_{(0)}y)$,
where $x_{(0)}y=[x,y]$ is the zeroth product (Lie bracket) extracted
from the simple-pole residue.
\end{computation}

\begin{remark}[Ordered vs.\ symmetric]
\label{rem:ordered-vs-symmetric-sl2}
In the symmetric bar complex, one would form
$\susp^{-1}e\otimes_\Sigma \susp^{-1}h$, identifying $e\otimes h$ with
$-h\otimes e$ (up to Koszul sign). In the ordered bar complex, both
live independently: the ordering of collisions on
$\operatorname{Conf}_2^{\mathrm{ord}}(\bC)$ is retained, not
merely whether a collision occurs. This is the surplus
information carried by the ordered bar complex of an
$E_\infty$-chiral algebra with poles.
\end{remark}

\begin{remark}[Duality interpretation]
The degree-$2$ bar cohomology, after linear dualisation, gives the
degree-$2$ component of the open-colour Koszul dual
$\mathcal A^!_{\mathrm{line}}=Y_\hbar(\mathfrak{sl}_2)$. The
relations~\eqref{eq:d-ef}--\eqref{eq:d-hh} dualise to the leading-order
Yangian commutation relations:
\[
[H^{(0)},E^{(0)}]=2E^{(0)},\qquad
[H^{(0)},F^{(0)}]=-2F^{(0)},\qquad
[E^{(0)},F^{(0)}]=H^{(0)}.
\]
These are exactly the $\mathfrak{sl}_2$ Lie bracket relations, confirming
that the leading-order piece of $Y_\hbar(\mathfrak{sl}_2)$ is
$U(\mathfrak{sl}_2)$.
\end{remark}

\subsubsection*{Degree 3: $d^2=0$ and the Yang--Baxter equation}

\begin{proposition}[$d^2=0$ is the classical Yang--Baxter equation;
\ClaimStatusProvedHere]
\label{prop:ybe-from-d-squared}
The condition $d^2=0$ on the degree-$3$ ordered bar complex
$\Barch^{\mathrm{ord}}_3(\widehat{\mathfrak{sl}}_2)$ is equivalent to
the classical Yang--Baxter equation
\begin{equation}\label{eq:cybe-sl2}
[r_{12}(z_1{-}z_2),\,r_{13}(z_1{-}z_3)]
\;+\;
[r_{12}(z_1{-}z_2),\,r_{23}(z_2{-}z_3)]
\;+\;
[r_{13}(z_1{-}z_3),\,r_{23}(z_2{-}z_3)]
\;=\;0
\end{equation}
for the collision residue $r(z)=\hbar\,\Omega/z$.
\end{proposition}

\begin{proof}
A degree-$3$ element $\susp^{-1}x\otimes \susp^{-1}y\otimes \susp^{-1}z$ has
two successive bar differentials. The first ($d_1$) acts on the
first two factors; the second ($d_2$) acts on the last two factors (with
a sign from the grading). Explicitly:
\[
d(\susp^{-1}x\otimes \susp^{-1}y\otimes \susp^{-1}z)
\;=\;
-\susp^{-1}[x,y]\otimes \susp^{-1}z
\;+\;
\susp^{-1}x\otimes \susp^{-1}[y,z].
\]
(The sign $+$ on the second term arises from $(-1)^{|\susp^{-1}x|}$.)
The full bar differential on three-point ordered configurations
incorporates the monodromy of the flat connection
$\nabla=d-r(z_{ij})\,dz_{ij}$ on
$\operatorname{Conf}_3^{\mathrm{ord}}(\bC)$, and the condition
$d^2=0$ for this corrected differential is precisely the CYBE~\eqref{eq:cybe-sl2}.

For $r(z)=\hbar\,\Omega/z$, the three commutator terms are:
\begin{align*}
[r_{12}(z_{12}),r_{13}(z_{13})]
&\;=\;\frac{\hbar^2}{z_{12}\,z_{13}}[\Omega_{12},\Omega_{13}],\\
[r_{12}(z_{12}),r_{23}(z_{23})]
&\;=\;\frac{\hbar^2}{z_{12}\,z_{23}}[\Omega_{12},\Omega_{23}],\\
[r_{13}(z_{13}),r_{23}(z_{23})]
&\;=\;\frac{\hbar^2}{z_{13}\,z_{23}}[\Omega_{13},\Omega_{23}].
\end{align*}
The partial-fraction identity
\[
\frac{1}{z_{12}\,z_{13}}\;=\;\frac{1}{z_{12}\,z_{23}}+\frac{1}{z_{13}\,z_{23}}
\qquad\text{(since }z_{13}=z_{12}+z_{23}\text{)}
\]
reduces the three-term sum to
$\frac{\hbar^2}{z_{12}\,z_{23}}
\bigl([\Omega_{12},\Omega_{13}]+[\Omega_{12},\Omega_{23}]\bigr)
+\frac{\hbar^2}{z_{13}\,z_{23}}
\bigl([\Omega_{13},\Omega_{12}]+[\Omega_{13},\Omega_{23}]\bigr)$.
Applying the Casimir identity
$[\Omega_{12},\Omega_{13}]+[\Omega_{12},\Omega_{23}]+[\Omega_{13},\Omega_{23}]=0$
(the Jacobi identity for~$\mathfrak g$ in tensor notation),
both grouped terms vanish, confirming the CYBE.
\end{proof}

\subsubsection*{The Yang $R$-matrix}

\begin{theorem}[Yang $R$-matrix from the ordered bar complex;
\ClaimStatusProvedHere]
\label{thm:yang-r-matrix}
The $R$-matrix on the fundamental representation $V=\bC^2$ of
$\mathfrak{sl}_2$, obtained from the ordered bar complex monodromy
\textup{(}the parallel transport of the connection
$\nabla=d-r^{\mathrm{coll}}(z)\,dz$ on
$\operatorname{Conf}_2^{\mathrm{ord}}(\bC)$\textup{)}, is the
\textbf{Yang $R$-matrix}:
\begin{equation}\label{eq:yang-r-matrix}
R(u)\;=\;u\,I\;+\;\hbar\,P
\;\in\;\operatorname{End}(V\otimes V)[u],
\end{equation}
where $P$ is the permutation operator ($P(v\otimes w)=w\otimes v$),
$u$ is the spectral parameter \textup{(}difference of evaluation
points\textup{)}, and $\hbar=1/(k+2)$.
\end{theorem}

\begin{proof}
The Casimir $\Omega$ acts on $V\otimes V\cong\bC^2\otimes\bC^2$
via $\rho\otimes\rho$, where $\rho:\mathfrak{sl}_2\to\mathfrak{gl}_2$ is
the fundamental representation. In the standard basis
$\{v_+,v_-\}$ of $V$:
\begin{align*}
\rho(e)&=E_{12}=\bigl(\begin{smallmatrix}0&1\\0&0\end{smallmatrix}\bigr),\quad
\rho(f)=E_{21}=\bigl(\begin{smallmatrix}0&0\\1&0\end{smallmatrix}\bigr),\quad
\rho(h)=\bigl(\begin{smallmatrix}1&0\\0&-1\end{smallmatrix}\bigr).
\end{align*}
The Casimir in the fundamental:
\[
(\rho\otimes\rho)(\Omega)
\;=\;E_{12}\otimes E_{21}+E_{21}\otimes E_{12}+\tfrac12(E_{11}-E_{22})\otimes(E_{11}-E_{22}).
\]
A direct matrix computation shows that
$(\rho\otimes\rho)(\Omega)=P-\tfrac12 I$, where $P$ is the
$4\times 4$ permutation matrix. (This is the standard identity for
$\mathfrak{sl}_2$: the Casimir in the tensor square of the
fundamental equals the flip minus half the identity.)

The $R$-matrix is read off from the leading collision residue as
\[
R(u)\;=\;u\,I+\hbar\,(\rho\otimes\rho)(\Omega)
\;=\;(u-\tfrac\hbar 2)\,I+\hbar\,P.
\]
After the standard shift $u\mapsto u+\hbar/2$ (which does not
affect the RTT relation), this becomes $R(u)=uI+\hbar P$, the Yang
$R$-matrix.

\textbf{Eigenvalues.} The space $V\otimes V$ decomposes as
$\operatorname{Sym}^2 V\oplus\bigwedge^2 V$ with dimensions $3+1$.
The permutation $P$ has eigenvalue $+1$ on $\operatorname{Sym}^2 V$
and $-1$ on $\bigwedge^2 V$. Therefore:
\begin{align}
R(u)\big|_{\operatorname{Sym}^2 V}&\;=\;(u+\hbar)\,\id,
\label{eq:r-eigenvalue-sym}\\
R(u)\big|_{\bigwedge^2 V}&\;=\;(u-\hbar)\,\id.
\label{eq:r-eigenvalue-wedge}
\end{align}
\end{proof}

\begin{remark}[Yang--Baxter verification]
The Yang $R$-matrix $R(u)=uI+\hbar P$ satisfies the quantum Yang--Baxter
equation
\[
R_{12}(u{-}v)\,R_{13}(u)\,R_{23}(v)
\;=\;R_{23}(v)\,R_{13}(u)\,R_{12}(u{-}v)
\]
in $\operatorname{End}(V^{\otimes 3})[u,v]$. This is the quantization
of the CYBE of Proposition~\ref{prop:ybe-from-d-squared}, confirming
that the degree-$3$ bar condition $d^2=0$ lifts to the full quantum
consistency condition on the ordered bar complex.
\end{remark}

\subsubsection*{The RTT presentation}

\begin{theorem}[RTT presentation of $Y_\hbar(\mathfrak{sl}_2)$;
\ClaimStatusProvedHere]
\label{thm:rtt-sl2}
Define the transfer matrix
\begin{equation}\label{eq:transfer-matrix-sl2}
T(u)\;=\;
\begin{pmatrix}
A(u)&B(u)\\C(u)&D(u)
\end{pmatrix}
\;\in\;\operatorname{End}(V)\otimes Y_\hbar(\mathfrak{sl}_2)[\![u^{-1}]\!],
\end{equation}
where $A(u),B(u),C(u),D(u)$ are formal power series in $u^{-1}$ with
coefficients in $Y_\hbar(\mathfrak{sl}_2)$. The RTT relation is
\begin{equation}\label{eq:rtt-sl2}
R_{12}(u{-}v)\,T_1(u)\,T_2(v)
\;=\;
T_2(v)\,T_1(u)\,R_{12}(u{-}v),
\end{equation}
where $R_{12}(u{-}v)=(u{-}v)I+\hbar P$ is the Yang $R$-matrix and
subscripts indicate the tensor factor in
$\operatorname{End}(V)\otimes\operatorname{End}(V)\otimes Y_\hbar$.

\textbf{The 16 component equations.}
Writing the RTT relation~\eqref{eq:rtt-sl2} in the standard basis
$\{v_+\otimes v_+,\,v_+\otimes v_-,\,v_-\otimes v_+,\,v_-\otimes v_-\}$
of $V\otimes V$ and comparing matrix entries, one obtains the following
generating-function relations (where we abbreviate $w=u-v$):
\begin{enumerate}[label=\textup{(\arabic*)}]
\item\label{rtt:AA}
$[A(u),A(v)]=0$,
\item\label{rtt:DD}
$[D(u),D(v)]=0$,
\item\label{rtt:BB}
$[B(u),B(v)]=0$,
\item\label{rtt:CC}
$[C(u),C(v)]=0$,
\item\label{rtt:AB}
$w\,A(u)B(v)=w\,B(v)A(u)+\hbar\,B(u)A(v)-\hbar\,A(v)B(u)$,
\item\label{rtt:AC}
$w\,A(u)C(v)=w\,C(v)A(u)-\hbar\,C(u)A(v)+\hbar\,A(v)C(u)$,
\item\label{rtt:DB}
$w\,D(u)B(v)=w\,B(v)D(u)-\hbar\,B(u)D(v)+\hbar\,D(v)B(u)$,
\item\label{rtt:DC}
$w\,D(u)C(v)=w\,C(v)D(u)+\hbar\,C(u)D(v)-\hbar\,D(v)C(u)$,
\item\label{rtt:BC}
$w\,[B(u),C(v)]=\hbar\bigl(D(u)A(v)-D(v)A(u)\bigr)$,
\item\label{rtt:AD}
$w\,[A(u),D(v)]=\hbar\bigl(C(u)B(v)-C(v)B(u)\bigr)$.
\end{enumerate}
Relations~\ref{rtt:AA}--\ref{rtt:CC} assert that the diagonal entries
commute among themselves and that $B$'s commute with $B$'s and $C$'s
with $C$'s. The content is in relations~\ref{rtt:AB}--\ref{rtt:AD},
which encode the Yangian deformation of $U(\mathfrak{sl}_2)$.
The remaining six of the sixteen matrix-entry equations are obtained by
applying the antiautomorphism $\sigma:A\leftrightarrow D$,
$B\leftrightarrow C$.
\end{theorem}

\begin{proof}
Direct expansion of~\eqref{eq:rtt-sl2} using $R(w)=wI+\hbar P$. The
left-hand side is $(wI+\hbar P)(T\otimes 1)(1\otimes T)$ and the
right-hand side is $(1\otimes T)(T\otimes 1)(wI+\hbar P)$; equating
the $4\times 4$ matrix entries produces the sixteen equations, of
which the listed ones are independent.
\end{proof}

\subsubsection*{Drinfeld generators via Gauss decomposition}

\begin{construction}[Gauss decomposition and Drinfeld generators;
\ClaimStatusProvedHere]
\label{constr:gauss-sl2}
The transfer matrix admits a Gauss decomposition:
\begin{equation}\label{eq:gauss-sl2}
T(u)\;=\;
\begin{pmatrix}1&0\\\mathcal{F}(u)&1\end{pmatrix}
\begin{pmatrix}H^+(u)&0\\0&H^-(u)\end{pmatrix}
\begin{pmatrix}1&\mathcal{E}(u)\\0&1\end{pmatrix},
\end{equation}
so that
\begin{align*}
A(u)&=H^+(u),&
B(u)&=H^+(u)\mathcal{E}(u),\\
C(u)&=\mathcal{F}(u)H^+(u),&
D(u)&=\mathcal{F}(u)H^+(u)\mathcal{E}(u)+H^-(u).
\end{align*}
The Drinfeld generators are the mode expansions:
\begin{align}
\mathcal{E}(u)&=\sum_{r\ge 0}\mathcal{E}^{(r)}\,u^{-r-1},\qquad
\mathcal{F}(u)=\sum_{r\ge 0}\mathcal{F}^{(r)}\,u^{-r-1},\label{eq:drinfeld-EF-sl2}\\
H^\pm(u)&=1+\sum_{r\ge 0}H^{\pm(r)}\,u^{-r-1}.\label{eq:drinfeld-H-sl2}
\end{align}
The leading modes $\mathcal{E}^{(0)}=e$, $\mathcal{F}^{(0)}=f$,
$H^{+(0)}-H^{-(0)}=h$ generate a copy of $U(\mathfrak{sl}_2)$ inside
$Y_\hbar(\mathfrak{sl}_2)$.
\end{construction}

The Drinfeld relations in these generators are:
\begin{align}
[H^+(u),H^+(v)]&=0,\qquad [H^-(u),H^-(v)]=0,
\label{eq:drinfeld-HH-sl2}\\
(u-v-\hbar)\,H^+(u)\,\mathcal{E}(v)&=(u-v+\hbar)\,\mathcal{E}(v)\,H^+(u),
\label{eq:drinfeld-HE-sl2}\\
(u-v+\hbar)\,H^-(u)\,\mathcal{E}(v)&=(u-v-\hbar)\,\mathcal{E}(v)\,H^-(u),
\label{eq:drinfeld-HmE-sl2}\\
[\mathcal{E}(u),\mathcal{F}(v)]
&=\frac{1}{\hbar(u-v)}\bigl(H^+(u)H^-(v)-H^+(v)H^-(u)\bigr),
\label{eq:drinfeld-EF-rel-sl2}\\
[\mathcal{E}(u),\mathcal{E}(v)]&=0,\qquad [\mathcal{F}(u),\mathcal{F}(v)]=0.
\label{eq:drinfeld-serre-sl2}
\end{align}
The last line is the $\mathfrak{sl}_2$ Serre relation (trivially
satisfied since $\mathfrak{sl}_2$ has a single simple root). For
$\mathfrak{sl}_N$ with $N\ge 3$, the Serre relations become the
non-trivial quantum Serre relations, arising from degree-$3$ bar
cohomology.

\subsubsection*{PBW basis and Hilbert series}

\begin{theorem}[PBW basis of $Y_\hbar(\mathfrak{sl}_2)$;
\ClaimStatusProvedHere]
\label{thm:pbw-sl2}
The Yangian $Y_\hbar(\mathfrak{sl}_2)$ admits a PBW basis:
ordered monomials in the generators
$\{\mathcal{E}^{(r)},\,\mathcal{F}^{(r)},\,H^{(r)}\}_{r\ge 0}$
\textup{(}in any fixed total order compatible with the mode
number~$r$\textup{)} form a basis of
$Y_\hbar(\mathfrak{sl}_2)$ as a filtered vector space. Consequently
the associated graded is a polynomial algebra in three generators per
mode:
\begin{equation}\label{eq:pbw-iso-sl2}
\gr\,Y_\hbar(\mathfrak{sl}_2)
\;\cong\;
U(\mathfrak{sl}_2[t])
\;\cong\;
\bC[\mathcal{E}^{(r)},\mathcal{F}^{(r)},H^{(r)}:\,r\ge 0].
\end{equation}
\end{theorem}

\begin{proof}
The PBW property for $Y_\hbar(\mathfrak g)$ is classical
(Drinfeld; see also Chari--Pressley).
For $\mathfrak g=\mathfrak{sl}_2$, it follows from the RTT
presentation: the RTT relations provide a complete rewriting system
that reduces any monomial in $\{A^{(r)},B^{(r)},C^{(r)},D^{(r)}\}$
to a linear combination of ordered monomials. The Gauss decomposition
translates this into a PBW basis for the Drinfeld generators.

From the bar complex perspective, the PBW property reflects the
\emph{Koszulness} of $\widehat{\mathfrak{sl}}_2$:
the ordered bar complex is formal (all higher $\Ainf$ operations
$m_k=0$ for $k\ge 3$), so the open-colour Koszul dual is quadratic
and PBW-complete. This is the pole-order class~$\mathbf L$
(double poles with finite shadow depth).
\end{proof}

\begin{corollary}[Hilbert series; \ClaimStatusProvedHere]
\label{cor:hilbert-sl2}
The Hilbert series of $Y_\hbar(\mathfrak{sl}_2)$, graded by mode
number, is
\begin{equation}\label{eq:hilbert-sl2}
\sum_{n\ge 0}\dim\bigl(Y_\hbar(\mathfrak{sl}_2)_{\le n}\bigr)\,q^n
\;=\;
\prod_{d\ge 1}\frac{1}{(1-q^d)^3}\,,
\end{equation}
where $Y_\hbar(\mathfrak{sl}_2)_{\le n}$ denotes the subspace spanned
by monomials of total mode number $\le n$, and the exponent
$3=\dim\mathfrak{sl}_2$ reflects the three Drinfeld generators
$\mathcal{E}^{(r)},\mathcal{F}^{(r)},H^{(r)}$ at each mode level.
\end{corollary}

\begin{proof}
Immediate from the PBW isomorphism~\eqref{eq:pbw-iso-sl2}: the
associated graded is a polynomial algebra with three generators at
each non-negative integral mode, giving the generating function
$\prod_{d\ge 1}(1-q^d)^{-3}$.
\end{proof}

\subsubsection*{Evaluation modules}

\begin{construction}[Evaluation modules; \ClaimStatusProvedHere]
\label{constr:eval-sl2}
For each $a\in\bC$, the \textbf{evaluation module}
$V_a=\bC^2$ is the $Y_\hbar(\mathfrak{sl}_2)$-module obtained
by evaluating the transfer matrix at a point:
\begin{equation}\label{eq:eval-module-sl2}
T(u)\;\mapsto\;I+\frac{\hbar\,(\rho\otimes\rho)(\Omega)}{u-a}
\;=\;I+\frac{\hbar\,P}{u-a}-\frac{\hbar}{2(u-a)}\,I
\;\in\;\operatorname{End}(V).
\end{equation}
Equivalently, in Drinfeld generators:
\begin{equation}\label{eq:eval-drinfeld-sl2}
\mathcal{E}^{(r)}\;\mapsto\;a^r\,\rho(e),\qquad
\mathcal{F}^{(r)}\;\mapsto\;a^r\,\rho(f),\qquad
H^{(r)}\;\mapsto\;a^r\,\rho(h).
\end{equation}
The parameter $a\in\bC$ is the \textbf{evaluation point}: it
is the position on the topological line $\bR_t\subset\bC_z\times\bR_t$
at which a boundary operator is inserted.
\end{construction}

\begin{proposition}[Tensor products and Yang--Baxter;
\ClaimStatusProvedHere]
\label{prop:eval-tensor-sl2}
The tensor product of evaluation modules $V_a\otimes V_b$ ($a\neq b$) is
irreducible as a $Y_\hbar(\mathfrak{sl}_2)$-module. The
$R$-matrix intertwining $V_a\otimes V_b\cong V_b\otimes V_a$ is
\begin{equation}\label{eq:eval-r-sl2}
\check R_{ab}\;=\;(a-b)\,I+\hbar\,P\;=\;R(a-b)
\;\in\;\operatorname{End}(V_a\otimes V_b),
\end{equation}
recovering the Yang $R$-matrix of Theorem~\ref{thm:yang-r-matrix}. At
$a=b$, the tensor product is reducible:
$V_a\otimes V_a\cong \operatorname{Sym}^2 V\oplus\bigwedge^2 V$ decomposes
as a direct sum of the $3$-dimensional irreducible and the trivial
representation. The $R$-matrix $R(0)=\hbar P$ is the
permutation, consistent with $R(0)$ having eigenvalue
$+\hbar$ on $\operatorname{Sym}^2$ and $-\hbar$ on $\bigwedge^2$.
\end{proposition}

\begin{proof}
The representation $V_a\otimes V_b$ is defined by the coproduct
$\Delta(T(u))=T(u)\dot\otimes T(u)$ applied to the evaluation
homomorphisms at $a$ and~$b$ respectively. The $R$-matrix
intertwiner is the Yang $R$-matrix evaluated at $u=a-b$, which is
invertible for $a\neq b$ (and has rank-$1$ kernel $\bigwedge^2 V$ at
$a-b=\hbar$). Irreducibility for generic $a-b$ follows from the
standard Drinfeld criterion.
\end{proof}

\subsubsection*{The open-colour Koszul dual identification}

\begin{theorem}[Open-colour Koszul dual of $\widehat{\mathfrak{sl}}_2$
is $Y_\hbar(\mathfrak{sl}_2)$; \ClaimStatusProvedHere]
\label{thm:sl2-koszul-dual}
Let $\mathcal A=\widehat{\mathfrak{sl}}_2$ at level~$k\neq -2$. The
open-colour Koszul dual
\[
\mathcal A^!_{\mathrm{line}}
\;=\;
H^*\bigl(\Barch^{\mathrm{ord}}(\mathcal A)\bigr)^\vee
\]
is isomorphic, as a filtered associative algebra with spectral parameter,
to the Yangian $Y_\hbar(\mathfrak{sl}_2)$ at deformation parameter
$\hbar=1/(k+2)$:
\begin{equation}\label{eq:koszul-dual-id-sl2}
\boxed{
\mathcal A^!_{\mathrm{line}}
\;\cong\;
Y_\hbar(\mathfrak{sl}_2),
\qquad
\hbar=\frac{1}{k+2}.
}
\end{equation}
\end{theorem}

\begin{proof}
Assemble the preceding computations:
\begin{enumerate}[label=(\roman*)]
\item \textbf{Generators.}
The degree-$1$ bar cohomology $H^1(\Barch^{\mathrm{ord}})$ is
$3$-dimensional, spanned by $[\susp^{-1}e]$, $[\susp^{-1}f]$,
$[\susp^{-1}h]$. Under linear dualisation, these become
the degree-$1$ generators of $\mathcal A^!_{\mathrm{line}}$,
identified with the Drinfeld currents $\mathcal E(u)$, $\mathcal F(u)$, $H(u)$.

\item \textbf{Relations.}
The degree-$2$ bar differential
(Computation~\ref{comp:ordered-bar-sl2})
dualises to the Yangian commutation relations at leading order.
The collision residue $r(z)=\hbar\Omega/z$ produces the spectral
deformation, upgrading $U(\mathfrak{sl}_2)$ relations to
$Y_\hbar(\mathfrak{sl}_2)$ relations.

\item \textbf{Consistency.}
The condition $d^2=0$ at degree~$3$
(Proposition~\ref{prop:ybe-from-d-squared}) is the CYBE, which
quantizes to the YBE satisfied by the Yang $R$-matrix
(Theorem~\ref{thm:yang-r-matrix}). This ensures the RTT presentation
(Theorem~\ref{thm:rtt-sl2}) is consistent.

\item \textbf{PBW completeness.}
The PBW basis (Theorem~\ref{thm:pbw-sl2}) confirms that no further
relations arise: the ordered bar complex is formal ($m_k=0$ for
$k\ge 3$), so the open-colour dual is quadratic and the PBW theorem
applies.

\item \textbf{Deformation parameter.}
The parameter $\hbar=1/(k+2)$ arises from
the collision residue~\eqref{eq:sl2-collision-residue}: the $d\log$
absorption of the double-pole OPE produces $r(z)=\hbar\Omega/z$,
and the bar differential encodes $\hbar$ as the deformation parameter
of the Yangian.
\end{enumerate}
The identification~\eqref{eq:koszul-dual-id-sl2} is thereby established.
\end{proof}

\begin{remark}[Comparison with the closed-colour dual]
\label{rem:closed-vs-open-sl2}
The closed-colour (symmetric, Francis--Gaitsgory) Koszul dual of
$\widehat{\mathfrak{sl}}_2$ at level~$k$ is
$\mathcal A^!_{\mathrm{ch}}=\widehat{\mathfrak{sl}}_2$ at level
$-k-2h^\vee=-k-4$ (Verdier-level duality). The open-colour
Koszul dual is the \emph{finite-type} Yangian $Y_\hbar(\mathfrak{sl}_2)$.
The input is the affine algebra; the closed output is another
affine algebra (at dual level); the open output is a finite-type
quantum group with spectral parameter. The cross-colour datum is the
$R$-matrix: $\Barch(\mathcal A)_n\simeq
(\Barch^{\mathrm{ord}}(\mathcal A)_n)^{R\text{-}\Sigma_n}$
(Proposition~\ref{prop:r-matrix-descent-vol1}).
\end{remark}

\begin{remark}[The $E_\infty/E_1$ interface]
\label{rem:e-inf-e1-sl2}
The input $\widehat{\mathfrak{sl}}_2$ is an $E_\infty$-chiral algebra:
it is a local vertex algebra whose chiral operations are defined on
the unordered configuration spaces $\operatorname{Conf}_n(X)$ with
$\Sigma_n$-equivariant factorisation structure. Its OPE has poles
(double poles from the level, simple poles from the Lie bracket), but
these are compatible with $E_\infty$; they are singularities of a
\emph{local} product.

The output $Y_\hbar(\mathfrak{sl}_2)$ is a genuinely $E_1$-chiral
algebra (in the sense of Etingof--Kazhdan): its $R$-matrix is
\emph{independent input}, not derivable from a local symmetric
structure. The passage from input to output (the
$E_1$-chiral Koszul duality functor) extracts the ordered bar
complex of an $E_\infty$ algebra and produces an $E_1$ algebra.
The $R$-matrix of the Yangian is independent of the input's
$E_\infty$ structure; it is the \emph{new} algebraic datum created by
the ordered bar construction.
\end{remark}

\begin{remark}[Formality: class~$\mathbf L$, shadow depth finite]
\label{rem:sl2-formality}
Like all affine Kac--Moody algebras, $\widehat{\mathfrak{sl}}_2$ belongs
to pole-order class~$\mathbf L$ (double poles, finite shadow depth).
The ordered bar complex is formal: all higher $\Ainf$ operations
$m_k=0$ for $k\ge 3$, and the open-colour Koszul dual $Y_\hbar(\mathfrak{sl}_2)$
is a quadratic algebra. This should be contrasted with the Virasoro
algebra (class~$\mathbf M$, quartic poles, infinite shadow depth),
whose open-colour Koszul dual carries genuinely infinite $\Ainf$
structure: the gravitational Yangian. The Drinfeld--Sokolov
reduction $\widehat{\mathfrak{sl}}_2\to\mathrm{Vir}_{c_{DS}}$
transports $\mathbf L\to\mathbf M$: the affine double-pole OPE
becomes the Virasoro quartic-pole OPE, and formality is destroyed.
\end{remark}
